\documentclass[12pt]{book}
\usepackage{amsmath}
\usepackage{amssymb}
\usepackage{geometry}
\geometry{margin=1in}

\title{Volume 13: Philosophy, Methodology, and Programmatics \\ Book 01: Methodology and Professionalization}
\author{James C. Harvey}
\date{December 2025}

\begin{document}
\maketitle
\tableofcontents

\chapter{Methodology as Geometry}

\section*{Abstract}
SDT methodology is not an external philosophy; it is the consequence of geometric derivation. This book
describes the programmatic structure of SDT as a discipline of consistent definitions, dimensional checks,
and derivation chains.

\section{Professionalization}
Professionalization means that every result is traceable to the master equation and its geometric
parameters. Documentation is therefore part of the physics, not merely its wrapper.

\chapter{Derivation Discipline}

\section*{Abstract}
Derivation discipline is the practice of tracing every claim to space, matter, movement, and now. This
chapter formalizes the derivation chain standard used across SDT.

\section{Chain Standard}
Every derivation must explicitly show:
\begin{equation}
\text{geometry} \rightarrow \text{occlusion} \rightarrow \text{pressure} \rightarrow \text{movement}.
\end{equation}

\chapter{Documentation as Formalism}

\section*{Abstract}
Documentation is a formal extension of the theory. It preserves the derivation chain and records the
conditions under which equations are valid.

\section{Documentation Rule}
Every document must define the boundary conditions, constants used, and the applicable regime.
\end{document}
